\section{Contrato y diseño propuesto}
\subsection{Semántica de la unidad aritmética}
Un registro de pesos \(W\) de 32 bits contiene ocho códigos U4. Un registro
\(A\) contiene cuatro activaciones S8. El selector \(h\in\{0,1\}\) elige
cuatro nibbles bajos o altos de \(W\), mientras \(z\in[0,15]\) es común a
los cuatro términos. Para \(i=0,\ldots,3\):
\begin{align}
q_{w,i}&=(W\gg4(i+4h))\mathbin{\&}15,\\
b_i&=(A\gg8i)\mathbin{\&}255,\\
q_{a,i}&=\begin{cases}b_i,&b_i<128,\\b_i-256,&b_i\ge128,\end{cases}\\
d&=\sum_{i=0}^{3}(q_{w,i}-z)q_{a,i}.
\end{align}
La salida es la representación INT32 de \(d\), sin saturación, escalas, bias
ni acumulador como tercer operando. El índice cero ocupa los bits bajos.
La elección lo/hi es independiente del orden de bytes en memoria; los programas
del proyecto usarán little-endian documentado.

\begin{evidencebox}{Ejemplo de contrato, no medición}
Con \(W=\code{0x00004A95}\), \(A=\code{0x0102FD04}\), \(z=8\) y \(h=0\),
los pesos corregidos son \([-3,1,2,-4]\) y las activaciones \([4,-3,2,1]\).
El subtotal es \(-15\), representado por \code{0xFFFFFFF1}. Es un ejemplo
matemático de la guía previa, no evidencia de una unidad RTL implementada.
\end{evidencebox}

\subsection{Rangos y diseño directo}
La resta de dos U4 abarca \([-15,15]\): requiere cinco bits con signo si
\(z\) es arbitrario. Un producto tiene valor absoluto máximo 1920 y cuatro
términos tienen cota 7680. Catorce bits con signo bastan para el subtotal;
se permite un árbol conservador de mayor anchura y extensión final a 32 bits.

\begin{center}
\begin{tikzpicture}[node distance=5mm]
\node[block,text width=28mm] (pick) {Seleccionar\\cuatro U4};
\node[block,fill=nord13,text width=28mm,right=of pick] (center) {Centrar\\U4 $-$ U4};
\node[block,fill=nord8,text width=28mm,right=of center] (mul) {Productos\\S5 $\times$ S8};
\node[block,fill=nord14,text width=28mm,right=of mul] (sum) {Reducir y\\extender signo};
\draw[arrow] (pick)--(center);\draw[arrow] (center)--(mul);\draw[arrow] (mul)--(sum);
\node[note,below=8mm of center,text width=125mm,xshift=18mm] {Ruta directa conceptual; el almacenamiento W4 no obliga a un multiplicador físico de 4 bits.};
\end{tikzpicture}
\captionof{figure}{Operaciones del datapath inicial. Las anchuras útiles deben verificarse antes de optimizar recursos.}
\end{center}
Con \(z=8\), el peso corregido sí cabe en S4. Esa especialización y la
factorización de la ecuación~\eqref{eq:factored} son alternativas legítimas;
no se descartan para favorecer la propuesta.

\subsection{Interfaz ISA abierta y control}
El nombre XQDot4Z y las variantes lo/hi son provisionales. Una instrucción
de 32 bits en espacio personalizado es una opción; el encoding no se congela
en esta versión. En particular, \(z\) será entrada de cuatro bits de la unidad,
pero su transporte desde el ISA sigue abierto.

Un inmediato solo representa un valor presente en el programa. Para zero-points
leídos en ejecución se necesitará despacho, especialización previa o una interfaz
adicional. Cualquier estado nuevo exige definir hazards, reset y cancelación;
su costo no se omite del experimento.

La primera unidad aislada será combinacional. Más adelante se podrán comparar
cuatro, dos y una vía. Número de pasos, latencia e intervalo de iniciación
son conceptos distintos. La integración deberá capturar operandos resueltos
por forwarding, retener una operación pendiente y transferir su resultado
una sola vez. Detener F/D/E mientras M/W drenan requiere burbujas o validez
equivalente; copiar repetidamente controles no es un protocolo correcto.

Se verificarán \(rd=x0\), fuentes/destino solapados, dependencias load-use,
consumo inmediato del resultado, stores, instrucciones consecutivas y flush/reset.
La corrección del RTL se comprobará contra una referencia entera independiente,
no contra la salida aproximada de una red neuronal.
